\section{Cadre conceptuel}
\label{sec:cadre_conceptuel}

Ce chapitre pose les concepts sur lesquels repose l'ensemble du travail. Il
définit d'abord les transferts de migrants et les notions qui s'y rattachent,
puis retrace les grandes approches de la pauvreté avant d'exposer celle qui est
retenue ici, la pauvreté multidimensionnelle, et la manière dont elle se compte.

\subsection{Définitions}

Les transferts de fonds des migrants sont définis, dans le cadre
statistique international du Fonds monétaire international, comme des
transferts personnels, c'est-à-dire l'ensemble des transferts courants, en
espèces ou en nature, effectués entre ménages résidents et ménages non
résidents, indépendamment du lien entre les personnes et du motif de
l'envoi \parencite{fmi2009}. Cette définition emporte trois traits qui
importent à l'analyse. Elle est d'abord large quant à la forme, couvrant
l'argent liquide remis de la main à la main, les envois par les canaux
formels, banques et sociétés de transfert, comme par les canaux informels,
ce qui correspond à la réalité sénégalaise où une part des fonds circule
hors du système bancaire. Elle est ensuite indifférente au motif, un envoi
pour la consommation courante, pour un évènement familial ou pour un
investissement relevant du même agrégat, alors que ces usages n'ont pas les
mêmes effets sur les conditions de vie. Elle est enfin établie au niveau du
ménage, et non de l'individu, ce qui rejoint la manière dont les enquêtes
la mesurent. La littérature du développement en retient la même substance,
les fonds envoyés par des personnes résidant à l'étranger vers leur pays ou
leur région d'origine \parencite{maimbo2005}, et en souligne le double
caractère, de revenu supplémentaire pour le ménage receveur et de lien
maintenu entre le migrant et sa famille \parencite{rapport2006}.

La migration désigne, selon le même glossaire, le déplacement de
personnes hors de leur lieu de résidence habituelle, soit à l'intérieur
d'un pays, soit au-delà d'une frontière internationale
\parencite{oim2019}. Elle se décline suivant plusieurs axes que la
littérature distingue soigneusement~: interne ou internationale selon que
la frontière est franchie~; temporaire, circulaire ou durable selon
l'horizon du déplacement~; de travail, d'études, familiale ou forcée selon
son motif. Les théories exposées dans la revue de littérature en rendent
compte à des niveaux différents, décision individuelle de placement de
capital humain \parencite{sjaastad1962}, arbitrage entre secteurs aux
revenus inégaux \parencite{harris1970}, stratégie familiale de
diversification du risque \parencite{stark1985}, et perpétuation par les
réseaux une fois le mouvement engagé \parencite{massey1993}. La migration
sénégalaise, ancienne et structurée par ces réseaux, relève pour l'essentiel
d'une migration internationale de travail, majoritairement masculine, dont
le ménage d'origine attend un soutien régulier.

Le migrant international est, selon le glossaire de l'Organisation
internationale pour les migrations, toute personne qui quitte son lieu de
résidence habituelle pour s'établir dans un autre pays, quelle que soit la
cause du déplacement, sa durée ou son caractère volontaire ou non
\parencite{oim2019}. La définition est volontairement neutre~: elle ne
distingue ni le travailleur du réfugié, ni l'installation durable du séjour
temporaire, distinctions qui relèvent de catégories juridiques propres à
chaque État. La littérature économique des transferts adopte un point de
vue complémentaire, celui du ménage d'origine~: le migrant y est un membre
parti dont le départ résulte le plus souvent d'une décision familiale, et
qui reste lié à ce ménage par des obligations réciproques dont l'envoi de
fonds est l'expression la plus courante \parencite{rapport2006,
stark1985}.

Le ménage est entendu au sens qu'en donne l'ANSD dans l'EHCVM, un groupe
de personnes, apparentées ou non, vivant sous le même toit, mettant en
commun leurs ressources et reconnaissant l'autorité d'un même chef de
ménage \parencite{ansd2019}. Cette définition, commune aux instituts de
statistique de l'UEMOA, repose sur la résidence et la communauté de
ressources plutôt que sur la parenté~: un ménage sénégalais peut réunir
plusieurs noyaux familiaux, et sa taille moyenne est parmi les plus élevées
de la sous-région, trait qui pèse sur la répartition interne des ressources
entre enfants. Le ménage bénéficiaire est celui qui reçoit un transfert de
migrant au cours d'une période de référence donnée~; le statut se définit à
ce niveau parce que les fonds entrent dans le budget commun, sans qu'aucune
information ne permette d'en suivre l'affectation à tel ou tel membre.

La pauvreté dispose elle aussi d'une définition internationalement
admise, arrêtée par le Sommet mondial pour le développement social de
Copenhague~: la pauvreté absolue est une situation caractérisée par une
privation grave des besoins humains fondamentaux, la nourriture, l'eau
potable, les installations sanitaires, la santé, le logement, l'éducation
et l'information, et elle dépend non seulement du revenu mais aussi de
l'accès aux services \parencite{onu1995}. Deux traits de cette définition
fondent la suite du travail~: elle est multidimensionnelle par construction,
la liste des besoins débordant ce que le revenu mesure, et elle désigne
l'accès aux services comme constitutif de la pauvreté, et non comme un
simple corrélat. La pauvreté des enfants en est la déclinaison propre à
l'enfance~: l'UNICEF la définit par la privation des ressources
matérielles, spirituelles et affectives nécessaires pour survivre, se
développer et s'épanouir, qui empêche les enfants de jouir de leurs droits
et d'atteindre leur plein potentiel \parencite{unicef2007, gordon2003}.

L'enfant, enfin, est toute personne âgée de moins de dix-huit ans, au
sens de l'article premier de la Convention internationale des droits de
l'enfant \parencite{onu1989}, définition que reprend la Charte africaine
des droits et du bien-être de l'enfant \parencite{oua1990}. Le seuil
d'âge n'est pas une commodité statistique~: il découle du régime de
protection que ces textes instituent, la personne mineure ne pouvant
subvenir seule à ses besoins ni faire valoir seule ses droits, si bien que
son bien-être dépend des ressources et des choix du ménage où elle vit.
C'est cette dépendance qui justifie de mesurer la pauvreté de l'enfant par
ses privations effectives plutôt que par le revenu du ménage.

L'évaluation d'impact, enfin, est définie par le manuel de référence de
la Banque mondiale comme l'analyse qui cherche à établir les changements
dans le bien-être des individus qui peuvent être attribués à un programme
ou à une politique, c'est-à-dire à répondre à une question de cause à
effet \parencite{gertler2016}. Sa difficulté constitutive est le
contrefactuel~: ce qu'il serait advenu des mêmes personnes sans
l'intervention ne s'observe jamais \parencite{heckman1997}. Toutes les
méthodes d'évaluation, expérimentation aléatoire, variables
instrumentales, régression sur discontinuité, appariement, double
différence, sont autant de manières de reconstruire ce contrefactuel à
partir de personnes non exposées, et se distinguent par les hypothèses
sous lesquelles cette reconstruction est valide
\parencite{angrist2009}. Le paramètre visé ici est l'effet moyen du
traitement sur les traités, l'écart entre ce que vivent les enfants des
ménages bénéficiaires et ce qu'ils auraient vécu sans les transferts~;
le choix de la méthode retenue et ses hypothèses sont exposés dans la
partie méthodologique.

\subsection{Approches de la pauvreté}

La pauvreté a reçu quatre grandes acceptions, dont chacune emporte une
mesure différente, et dont la succession dessine un élargissement
progressif : du revenu aux biens, des biens aux conditions de vie, des
conditions de vie aux capacités. Situer la mesure retenue ici suppose de
les parcourir dans cet ordre.

L'approche monétaire identifie la pauvreté à l'insuffisance du revenu
ou de la consommation au regard d'un seuil. Sa force est sa simplicité et sa
comparabilité~; sa faiblesse, pour la question posée ici, est double. Elle
suppose que les ressources se répartissent équitablement à l'intérieur du
ménage, ce que rien n'assure \parencite{deaton1997}, et elle ignore ce que le
revenu ne peut acheter faute d'une offre disponible, une école ou un réseau
d'assainissement.

L'approche par les besoins essentiels y répond en dressant la liste des
biens et services jugés indispensables, et en déclarant pauvre celui qui n'y
accède pas. C'est de cette tradition que procède la mesure des privations des
enfants développée par \textcite{gordon2003} pour l'UNICEF, qui fonde les
approches mobilisées dans ce travail. Elle déplace la mesure du revenu vers
les conditions de vie, mais laisse ouverte la question du choix des besoins
retenus.

L'approche relative de \textcite{townsend1979} définit la privation par
rapport aux conditions de vie jugées normales dans une société donnée : la
pauvreté n'y est plus un manque absolu, mais l'impossibilité de participer
à la vie ordinaire de son milieu. Son apport décisif est d'avoir déplacé la
mesure du revenu vers une liste d'items de conditions de vie observés
directement, logement, alimentation, équipement, relations sociales, dont
le cumul des manques définit la privation. C'est cette démarche par
indicateurs multiples que les mesures multidimensionnelles contemporaines
systématisent.

L'approche par les capabilités de \textcite{sen1985, sen1999} opère le
déplacement décisif. Ce qui importe n'est ni le revenu ni la possession de
biens, mais la capacité effective d'une personne à être et à faire ce qu'elle a
des raisons de valoriser. Le revenu n'est plus qu'un moyen, dont le rendement
en capacités varie d'un individu et d'un contexte à l'autre. C'est dans ce
cadre que s'inscrit la mesure retenue ici.

% ─────────────────────────────────────────────────────────────────────────────
\subsection{Pauvreté multidimensionnelle de l'enfant}
\label{subsubsec:concept_moda}
% ─────────────────────────────────────────────────────────────────────────────

La mesure monétaire de la pauvreté n'est pas adaptée aux enfants, elle ne capte
ni l'accès aux services publics, ni la répartition intra-ménage des ressources,
qu'elle suppose équitable à tort \parencite{deaton1997}, ni les besoins
spécifiques de l'enfant \parencite{gordon2003}. Le cadre des \og capabilités \fg{}
de \textcite{sen1999, sen1985} y substitue une définition en termes de privation
de capacités fondamentales, opérationnalisée par \textcite{townsend1979} puis
par \textcite{unicef2007}. Ce qui importe est la capacité effective de l'enfant
à fonctionner dans chaque dimension de son développement, d'autant que les
privations de l'enfance ont des effets durables sur le capital humain.

La pauvreté multidimensionnelle des enfants se définit dans ce cadre comme
la privation simultanée dans plusieurs dimensions constitutives du bien-être
de l'enfant. Sa mesure obéit à trois exigences que la littérature a
progressivement dégagées \parencite{roelen2008, deneubourg2012}. La première
est de prendre l'enfant individuel, et non le ménage, comme unité d'analyse,
car les indices hérités de la tradition centrée sur le ménage ne reflètent pas
directement les privations propres à chaque enfant, et deux enfants d'un
même ménage peuvent connaître des situations différentes. La deuxième est
d'adapter les indicateurs à l'âge, les besoins fondamentaux variant
considérablement entre un nourrisson et un adolescent, et appliquer les mêmes
indicateurs et les mêmes seuils à tous les âges revient à mal mesurer chacun
d'eux. La troisième est de
s'intéresser aux privations superposées, c'est-à-dire aux privations qui se
combinent chez les enfants les plus vulnérables et à leur intensité,
plutôt qu'à chaque manque pris isolément.

Cette conception s'ancre enfin dans les droits de l'enfant, tels que les
consacrent la Convention internationale des droits de l'enfant
\parencite{onu1989} et la Charte africaine des droits et du bien-être de
l'enfant \parencite{oua1990}, chaque dimension du bien-être correspondant à
un droit garanti par ces textes. La santé et la nutrition relèvent du droit
à la survie et au meilleur état de santé possible (article 24 de la convention,
article 14 de la charte)~; l'éducation, du droit à l'instruction
(articles 28 et 29 de la convention, article 11 de la charte)~; l'eau,
l'assainissement et le logement, du droit à un niveau de vie suffisant
(article 27 de la convention)~; la protection, du droit d'être préservé de toute
forme d'exploitation, de séparation abusive d'avec ses parents, et d'être
enregistré à la naissance (articles 7, 9 et 32 de la convention, articles 6 et 15
de la charte). Une privation n'est donc pas un simple manque de ressources
mais la négation d'un droit, ce qui fonde la mesure sur une définition
normative du bien-être de l'enfant plutôt que sur le seul revenu du ménage.

% ─────────────────────────────────────────────────────────────────────────────
\subsection{Mesure des privations}
% ─────────────────────────────────────────────────────────────────────────────

Définir la pauvreté comme une privation simultanée ne suffit pas, encore
faut-il la compter. La méthode de dénombrement de \textcite{alkire2011} fournit
le cadre aujourd'hui dominant. Elle procède en deux temps. Un premier seuil,
dit intra-dimensionnel, détermine pour chaque dimension si l'enfant y est privé.
Un second seuil, dit inter-dimensionnel et noté $k$, fixe le nombre de
dimensions à partir duquel l'enfant est déclaré pauvre. Ce double seuil évite
les deux écueils symétriques des approches antérieures, l'approche dite de
l'union qui déclare pauvre tout enfant privé dans une seule dimension et
surestime donc la pauvreté, et l'approche dite de l'intersection qui exige la
privation dans toutes les dimensions et la sous-estime tout autant.

De ce dénombrement se déduisent trois indices, qui ne disent pas la même chose.
L'incidence $H$ est la part des enfants privés dans au moins $k$ dimensions.
L'intensité $A$ est la part moyenne de privations que supportent les enfants
ainsi déclarés pauvres. L'indice ajusté $M_0$, produit des deux, est le seul
des trois à satisfaire la propriété de monotonicité dimensionnelle, une
privation supplémentaire chez un enfant déjà pauvre augmentant l'indice alors
qu'elle laisse l'incidence inchangée. Suivre les trois est donc nécessaire, une
baisse de $H$ accompagnée d'une intensité stable ne signifiant pas la même
chose qu'une baisse conjointe des deux.

L'approche MODA \parencite{deneubourg2012} applique ce cadre à l'enfant. Elle
en retient la logique du double seuil, mais l'assortit de trois choix propres
au public visé. Les dimensions sont celles du bien-être de l'enfant et non du
ménage, elles sont déclinées par groupe d'âge, et l'analyse porte sur les
privations qui se superposent chez un même enfant plutôt que sur les taux
dimension par dimension. C'est cette déclinaison, adaptée au Sénégal par
\textcite{ansd2024}, que ce mémoire reprend.

Deux limites de cette famille de mesures doivent être signalées, car elles
pèsent sur l'interprétation des résultats. La première tient à la pondération,
les dimensions étant ici traitées comme équivalentes faute d'un critère
consensuel pour les hiérarchiser, choix courant mais qui n'est pas neutre. La
seconde tient au caractère discret des seuils, une amélioration réelle mais
insuffisante pour franchir un seuil restant invisible à l'indice, ce qui rend
nécessaire de tester la sensibilité des résultats à la valeur de $k$.
